\clearpage
\begin{examplebox}
\textbf{例 1：最直观的整点验证}

已知 $C:y^2=4x$，$A=(1,2)$，$M=(0,0)$，$N=(4,-4)$。

\textbf{核对点：}$A,M,N$ 都在抛物线上，三点互不相同。
\[
 k_{AM}=\frac{0-2}{0-1}=2,\qquad
 k_{AN}=\frac{-4-2}{4-1}=-2.
\]
于是满足斜率对称条件，且
\[
 \boxed{k_{MN}=\frac{-4}{4}=-1},\qquad MN:\ y=-x.
\]
\textbf{提醒：}焦点为 $F=(1,0)$，不能把 $F$ 误写成 $A$。
\end{examplebox}

\clearpage
\begin{examplebox}
\textbf{例 2：两动点不关于 $x$ 轴对称，也照样成立}
\begin{center}\begin{tikzpicture}[font=\small,>=stealth,line join=round]
\draw[gray!45,rounded corners=3pt] (0,0) rectangle (14.2,7.2);
\node at (7.1,6.6) {分别找直线的第二交点，纵坐标中的变动部分相消};
\begin{scope}[shift={(1.4,3.8)},scale=.5]
\draw[->,gray] (-.5,0)--(9.7,0) node[above] {$x$};
\draw[->,gray] (0,-4.8)--(0,3.7) node[left] {$y$};
\draw[teal!65!black,thick,samples=110,domain=-4.3:3.4,variable=\t] plot ({\t*\t/2},\t);
\draw[gray,dashed] (-.2,1)--(4.5,1);
\draw[blue!75!black,very thick] (.5,1)--(2,2);
\draw[orange!85!black,very thick] (.5,1)--(8,-4);
\draw[green!45!black,very thick] (2,2)--(8,-4);
\fill (.5,1) circle(2pt) node[above left=10pt] {$A$};
\fill (2,2) circle(2pt) node[above=12pt] {$M$};
\fill (8,-4) circle(2pt) node[above right=12pt] {$N$};
\end{scope}
\node[anchor=west,align=left] at (8.8,3.8)
{$p=1,\quad t=\frac23$\\[14pt]
$u=\dfrac{2p}{t}-p=2$\\[16pt]
$v=-\dfrac{2p}{t}-p=-4$\\[16pt]
$u+v=-2p=-2$};
\node[text width=13.2cm,align=center] at (7.1,.65) {蓝、橙线分别与曲线在 $A$ 之外再交于 $M,N$；绿色弦斜率为 $-1$。};
\end{tikzpicture}
\end{center}


取 $C:y^2=2x$，$p=1$，$A=(\frac12,1)$；令
\[
 M=(2,2),\qquad N=(8,-4).
\]
计算：
\[
 k_{AM}=\frac{2-1}{2-\frac12}=\frac23,\qquad
 k_{AN}=\frac{-4-1}{8-\frac12}=-\frac23,
\]
\[
 k_{MN}=\frac{-4-2}{8-2}=\boxed{-1}.
\]
此时两点纵坐标 $2$ 与 $-4$ 并不互为相反数；正确的不变量是
\[
 2+(-4)=-2=-2p.
\]
\end{examplebox}

\clearpage
\begin{examplebox}
\textbf{例 3：从 $k_{MN}=-1$ 反过来判断斜率对称}

仍取 $C:y^2=4x$，$A=(1,2)$，但改为
\[
 M=\left(\frac14,-1\right),\qquad
 N=\left(\frac94,-3\right).
\]
有 $u+v=-1-3=-4=-2p$，因此 $k_{MN}=-1$。
\[
 k_{AM}=\frac{2p}{p+u}=\frac4{2-1}=4,\qquad
 k_{AN}=\frac4{2-3}=-4.
\]
这说明：\textbf{在三点互异且 $AM,AN$ 斜率有限的前提下，原结论还可反向使用：}
\[
 \boxed{k_{AM}=-k_{AN}\ \Longleftrightarrow\ k_{MN}=-1}.
\]
\end{examplebox}

\clearpage
\begin{examplebox}
\textbf{例 4：推广——定点改变时，第三条弦的斜率怎样变化？}
\begin{center}\begin{tikzpicture}[font=\small,>=stealth,line join=round]
\draw[gray!45,rounded corners=3pt] (0,0) rectangle (14.2,10.5);
\node at (7.1,9.9) {一般定点的对照：纵坐标为零时，第三弦改为竖直线};
\node at (3.55,8.5) {$a=p=1$};
\node at (10.65,8.5) {$a=0$};
\begin{scope}[shift={(2.4,4.9)},scale=1.0]
\draw[->,gray] (-.7,0)--(3.1,0) node[above] {$x$};
\draw[->,gray] (0,-2.6)--(0,2.7) node[left] {$y$};
\draw[teal!65!black,thick,samples=80,domain=-2.3:2.3,variable=\t] plot ({\t*\t/2},\t);
\draw[blue!75!black,thick] (.5,1)--(0,0);
\draw[orange!85!black,thick] (.5,1)--(2,-2);
\draw[green!45!black,very thick] (0,0)--(2,-2);
\fill (.5,1) circle(2pt) node[right=14pt] {$A_0$};
\fill (0,0) circle(2pt) node[below left=9pt] {$M$};
\fill (2,-2) circle(2pt) node[right=12pt] {$N$};
\end{scope}
\begin{scope}[shift={(9.5,4.9)},scale=1.0]
\draw[->,gray] (-.7,0)--(3.1,0) node[above] {$x$};
\draw[->,gray] (0,-2.6)--(0,2.7) node[left] {$y$};
\draw[teal!65!black,thick,samples=80,domain=-2.3:2.3,variable=\t] plot ({\t*\t/2},\t);
\draw[blue!75!black,thick] (0,0)--(.5,1);
\draw[orange!85!black,thick] (0,0)--(.5,-1);
\draw[green!45!black,very thick] (.5,1)--(.5,-1);
\fill (0,0) circle(2pt) node[below left=9pt] {$A_0=O$};
\fill (.5,1) circle(2pt) node[right=14pt] {$M$};
\fill (.5,-1) circle(2pt) node[right=14pt] {$N$};
\end{scope}
\node at (3.55,1.6) {$k_{MN}=-1$};
\node at (10.65,1.6) {$MN:x=\frac12$，斜率不存在};
\node[text width=13.2cm,align=center] at (7.1,.65) {两图同取 $p=1$、相同单位长度；$a=0$ 时不能代入 $-p/a$。};
\end{tikzpicture}
\end{center}


设定点改为 $A_0=(\frac{a^2}{2p},a)$，其中 $a\ne0$ 且仍在抛物线上；$M,N$ 互异且不同于 $A_0$，相关斜率存在。

若 $k_{A_0M}=-k_{A_0N}$，则
\[
 \frac{2p}{a+u}=-\frac{2p}{a+v}
 \quad\Longrightarrow\quad u+v=-2a.
\]
所以
\[
 \boxed{k_{MN}=\frac{2p}{u+v}=-\frac pa}.
\]
当 $a=p$ 时即回到 C098 的 $-1$；当 $a=0$ 时，推导给出 $u+v=0$，此时 $MN$ 为\textbf{竖直线}，不能套入 $-p/a$。
\end{examplebox}
